% Section 11.3, from lectures/L11/appendix/practice_test/questions.md; the rubric from the same
% file and from info/README.md.

\needspace{16\baselineskip}
\section{Övningsdugga 2}
\label{c11:sec:ovningsdugga}
Övningsduggan omfattar kapitel~\ref{c7:ch}--\ref{c10:ch}: vektorer, funktioner, trigonometri,
logaritmer och decibel. Den är upplagd som en riktig dugga, så skriv den på egen hand och i ett
svep innan du tittar i facit, där svaren finns.

\begin{aside}
\textbf{Så går duggan till.} Tillåtna hjälpmedel är skrivmaterial, valfri miniräknare och de
formler ur kursens formelsamling som behövs. Duggan har fem uppgifter om $1{,}0$ poäng var, alltså
högst $5{,}0$ poäng. Alla uppgifter kräver lösningar med redovisat svar, inklusive enhet där det
behövs. Den riktiga duggan har samma upplägg: den påverkar inte betyget, men minst $3{,}5$ av
$5{,}0$ poäng ger $0{,}5$ bonuspoäng på den ordinarie tentamen.
\end{aside}

\prov{od2}

\provuppgift{1}{1,0}
Du har två vektorer $\mathbf{u} = (1;\,2)$ och $\mathbf{v} = (-3;\,4)$.
\begin{parts}
  \item Rita upp vektorerna i ett koordinatsystem.
  \item Beräkna vektorernas absolutbelopp $\abs{\mathbf{u}}$ och $\abs{\mathbf{v}}$.
  \item Beräkna vektorernas vinklar $v_u$ och $v_v$.
  \item Beräkna en tredje vektor $\mathbf{w} = 2\mathbf{u} - \mathbf{v}$.
\end{parts}

\provuppgift{2}{1,0}
Förenkla följande rationella uttryck och ange vilka $x$-värden som uttrycket inte får anta:
\[
  f(x) = \frac{x^2 - x - 2}{x - 2}
\]

\provuppgift{3}{1,0}
När en kondensator urladdas genom ett motstånd minskar spänningen exponentiellt med tiden.
Spänningen $u(t)$ kan beskrivas med
\[
  u(t) = U_0 e^{-t/(RC)},
\]
där
\begin{itemize}
  \item $u(t)$ är spänningen över kondensatorn vid tiden $t$,
  \item $U_0$ är begynnelsespänningen i V,
  \item $RC$ är kretsens tidskonstant i sekunder,
  \item $t$ är tiden i sekunder.
\end{itemize}
En kondensator med kapacitansen $470\,\mu\text{F}$ är ansluten till ett motstånd på
$10\,\text{k}\Omega$. Den är initialt laddad till $12\,\text{V}$ och börjar urladdas vid $t = 0$.
\begin{parts}
  \item Beräkna spänningen efter fem sekunder.
  \item Bestäm funktionens definitionsmängd och värdemängd för tidsintervallet
    $0 \leq t \leq 20\,\text{s}$.
  \item Beräkna efter hur lång tid spänningen har sjunkit till hälften av sitt ursprungliga
    värde.
\end{parts}

\provuppgift{4}{1,0}
En växelspänning har amplituden $5\,\text{V}$, frekvensen $100\,\text{Hz}$ och fasen
$90^{\circ}$.
\begin{parts}
  \item Bestäm växelspänningens ekvation $u(t)$. Ange fasen i radianer.
  \item Rita växelspänningens sinuskurva över en period $T$.
\end{parts}

\provuppgift{5}{1,0}
En ljudförstärkare har en förstärkning på $32\,\text{dB}$. Hur många gångers
spänningsförstärkning motsvarar det?
